\noindent
\begin{minipage}[t]{0.28\textwidth}
    \strut
\end{minipage}%
\hfill
\begin{minipage}[t]{0.70\textwidth}
    Chúng ta tóm tắt các công thức đạo hàm đã học cho đến nay như sau:
    \vspace{4pt}

    \begin{tcolorbox}[
        colframe=stewartred,
        colback=white,
        boxrule=0.8pt,
        arc=0mm,
        top=5pt, bottom=6pt, left=8pt, right=8pt
    ]
        {\bfseries\sffamily\color{stewartred}Bảng các công thức đạo hàm}
        \vspace{2pt}
        \begin{center}
        \renewcommand{\arraystretch}{1.85}
        \begin{tabular*}{\linewidth}{@{\extracolsep{\fill}} lll @{}}
            $\dfrac{d}{dx}(c) = 0$ & $\dfrac{d}{dx}(x^n) = nx^{n-1}$ & $\dfrac{d}{dx}(e^x) = e^x$ \\[3pt]
            $(cf)' = cf'$ & $(f + g)' = f' + g'$ & $(f - g)' = f' - g'$ \\[3pt]
            $(fg)' = fg' + gf'$ & $\left(\dfrac{f}{g}\right)' = \dfrac{gf' - fg'}{g^2}$ & 
        \end{tabular*}
        \end{center}
    \end{tcolorbox}
\end{minipage}

\vspace{14pt}

\noindent
\begin{tikzpicture}
    \node[anchor=west, inner sep=0pt] (title) at (1.4cm, 0) {%
        {\fontsize{19}{22}\selectfont\bfseries\sffamily\color{stewartcyan}3.2}\hspace{7pt}%
        {\color{stewartcyan}\rule[-3pt]{1.5pt}{18pt}}\hspace{7pt}%
        {\fontsize{16}{19}\selectfont\bfseries\sffamily Bài tập}%
    };
    \foreach \y in {-3pt, -1pt, 1pt, 3pt} {
        \draw[stewartcyan, line width=0.6pt] (0, \y) -- ([xshift=-8pt]title.west |- 0, \y);
    }
    \foreach \y in {-3pt, -1pt, 1pt, 3pt} {
        \draw[stewartcyan, line width=0.6pt] ([xshift=8pt]title.east |- 0, \y) -- (\linewidth, \y);
    }
\end{tikzpicture}

\vspace{8pt}

\begin{multicols}{2}
\setlength{\columnsep}{22pt}

\noindent
\textbf{1.}\enspace Tìm đạo hàm của $f(x) = (1 + 2x^2)(x - x^2)$ bằng hai cách: sử dụng Quy tắc tích và thực hiện phép nhân trước. Các kết quả của bạn có khớp nhau không?

\vspace{5pt}
\noindent
\textbf{2.}\enspace Tìm đạo hàm của hàm số
\[
F(x) = \frac{x^4 - 5x^3 + \sqrt{x}}{x^2}
\]
bằng hai cách: sử dụng Quy tắc thương và rút gọn trước. Hãy chỉ ra rằng các kết quả của bạn là tương đương. Bạn thích phương pháp nào hơn?

\vspace{6pt}
\noindent
{\color{stewartcyan}\textbf{\textsf{3–30}}}\quad \textbf{Tính đạo hàm.}

\vspace{5pt}
\noindent
\textbf{3.}\enspace $y = (4x^2 + 3)(2x + 5)$ \par\vspace{5.5pt}

\noindent
\textbf{4.}\enspace $y = (10x^2 + 7x - 2)(2 - x^2)$ \par\vspace{5.5pt}

\twoex{5}{$y = x^3 e^x$}{6}{$y = (e^x + 2)(2e^x - 1)$}
\twoex{7}{$f(x) = (3x^2 - 5x)e^x$}{8}{$g(x) = (x + 2\sqrt{x})e^x$}
\twoex{9}{$y = \dfrac{x}{e^x}$}{10}{$y = \dfrac{e^x}{1 - e^x}$}
\twoex{11}{$g(t) = \dfrac{3 - 2t}{5t + 1}$}{12}{$G(u) = \dfrac{6u^4 - 5u}{u + 1}$}
\twoex{13}{$f(t) = \dfrac{5t}{t^3 - t - 1}$}{14}{$F(x) = \dfrac{1}{2x^3 - 6x^2 + 5}$}
\twoex{15}{$y = \dfrac{s - \sqrt{s}}{s^2}$}{16}{$y = \dfrac{\sqrt{x}}{\sqrt{x} + 1}$}

\noindent
\textbf{17.}\enspace $J(u) = \left(\dfrac{1}{u} + \dfrac{1}{u^2}\right)\left(u + \dfrac{1}{u}\right)$ \par\vspace{5.5pt}

\noindent
\textbf{18.}\enspace $h(w) = (w^2 + 3w)(w^{-1} - w^{-4})$ \par\vspace{5.5pt}

\noindent
\textbf{19.}\enspace $H(u) = (u - \sqrt{u})(u + \sqrt{u})$

\columnbreak

\noindent
\textbf{20.}\enspace $f(z) = (1 - e^z)(z + e^z)$ \par\vspace{5.5pt}

\twoex{21}{$V(t) = (t + 2e^t)\sqrt{t}$}{22}{$W(t) = e^t(1 + te^t)$}
\twoex{23}{$y = e^p(p + p\sqrt{p})$}{24}{$h(r) = \dfrac{ae^r}{b + e^r}$}
\twoex{25}{$f(t) = \dfrac{\sqrt[3]{t}}{t - 3}$}{26}{$y = (z^2 + e^z)\sqrt{z}$}
\twoex{27}{$f(x) = \dfrac{x^2 e^x}{x^2 + e^x}$}{28}{$F(t) = \dfrac{At}{Bt^2 + Ct^3}$}
\twoex{29}{$f(x) = \dfrac{x}{x + \dfrac{c}{x}}$}{30}{$f(x) = \dfrac{ax + b}{cx + d}$}

\vspace{7pt}
\noindent
{\color{stewartcyan}\textbf{\textsf{31–34}}}\quad \textbf{Tìm $f'(x)$ và $f''(x)$.}

\vspace{5pt}
\twoex{31}{$f(x) = x^2 e^x$}{32}{$f(x) = \sqrt{x}e^x$}
\twoex{33}{$f(x) = \dfrac{x}{x^2 - 1}$}{34}{$f(x) = \dfrac{x}{1 + \sqrt{x}}$}

\vspace{7pt}
\noindent
{\color{stewartcyan}\textbf{\textsf{35–36}}}\quad \textbf{Tìm phương trình tiếp tuyến của đường cong đã cho tại điểm được chỉ ra.}

\vspace{5pt}
\twoex{35}{$y = \dfrac{x^2}{1 + x},\quad \left(1, \dfrac{1}{2}\right)$}{36}{$y = \dfrac{1 + x}{1 + e^x},\quad \left(0, \dfrac{1}{2}\right)$}

\vspace{7pt}
\noindent
{\color{stewartcyan}\textbf{\textsf{37–38}}}\quad \textbf{Tìm các phương trình của tiếp tuyến và pháp tuyến của đường cong đã cho tại điểm được chỉ ra.}

\vspace{5pt}
\twoex{37}{$y = \dfrac{3x}{1 + 5x^2},\quad \left(1, \dfrac{1}{2}\right)$}{38}{$y = x + xe^x,\quad (0, 0)$}

\end{multicols}
